\subsection{Products of manifolds}

5.6.1. Let X and X' be two sets and let \(c = (\mathrm{U}, \varphi, \mathrm{E})\) (resp. \(c' = (\mathrm{U}', \varphi, \mathrm{E}')\) ) be a chart of X (resp. X'). The triple

\[
(\mathbf {U} \times \mathbf {U} ^ {\prime}, \varphi \times \varphi^ {\prime}, \mathbf {E} \times \mathbf {E} ^ {\prime})
\]

is a chart of the product set \(X \times X'\) , denoted \(c \times c'\) .

5.6.2. Let X and X' be two manifolds of class \(C^{r}\). On the product set X × X' there exists one and only one manifold structure of class \(C^{r}\) such that c × c' is a chart of X × X' for every chart c of X and every chart c' of X'. Equipped with this structure, X × X' is called the product manifold of the manifolds X and X'. The product of any finite number of manifolds is defined similarly. The topology of the manifold X × X' is the product topology of the topologies of X and X'. For a ∈ X and b ∈ X', we have:

\[
\dim_ {(a, b)} (\mathrm{X} \times \mathrm{X} ^ {\prime}) = \dim_ {a} \mathrm{X} + \dim_ {b} \mathrm{X} ^ {\prime}.
\]

5.6.3. Let X and X' be two manifolds, and let X × X' be their product. Let a ∈ X and let a' ∈ X'. The canonical projections

\[
\operatorname{pr} _ {1}: \mathrm{X} \times \mathrm{X} ^ {\prime} \rightarrow \mathrm{X} \quad \text { and } \quad \operatorname{pr} _ {2}: \mathrm{X} \times \mathrm{X} ^ {\prime} \rightarrow \mathrm{X} ^ {\prime}
\]

are morphisms. Let \(\pi_{i}\) ( \(i=1,2\) ) be their tangent linear maps at the point \((a,a')\) . The map

\[
(\pi_ {1}, \pi_ {2}): \mathrm{T} _ {(a, a ^ {\prime})} (\mathrm{X} \times \mathrm{X} ^ {\prime}) \rightarrow \mathrm{T} _ {a} (\mathrm{X}) \times \mathrm{T} _ {a ^ {\prime}} (\mathrm{X} ^ {\prime})
\]

is an isomorphism, which allows one to identify the tangent space to X × X' at (a, a') with the product \(T_{a}(X) \times T_{a'}(X')\).

The injection of \(\mathbf{T}_{a}(\mathbf{X})\) to \(\mathbf{T}_{(a,a^{\prime})}(\mathbf{X}\times \mathbf{X}^{\prime})\) resulting from this identification is the tangent linear map at \(a\) to the morphism \(x\mapsto (x,a^{\prime})\) of \(\mathbf{X}\) to \(\mathbf{X}\times \mathbf{X}^{\prime}\) ; we have an analogous result for the injection of \(\mathbf{T}_{a^{\prime}}(\mathbf{X}^{\prime})\) to \(\mathbf{T}_{(a,a^{\prime})}(\mathbf{X}\times \mathbf{X}^{\prime})\) .

5.6.4. Let W, X and X' be three manifolds, and let \(f: \mathbf{W} \to \mathbf{X}\) , \(f': \mathbf{W} \to \mathbf{X'}\) be two maps. For the map

\[
(f, f ^ {\prime}): \mathbf {W} \rightarrow \mathbf {X} \times \mathbf {X} ^ {\prime}
\]

to be a morphism, it is necessary and sufficient that \(f\) and \(f'\) be morphisms (this justifies the use of the term "product", cf.~Ens., chap.~IV, § 2, no. 4). In this case, if \(a\) is a point of W, we have:

\[
\mathrm{T} _ {a} (f, f ^ {\prime}) = (\mathrm{T} _ {a} (f), \mathrm{T} _ {a} (f ^ {\prime})),
\]

taking into account the identification made above.

5.6.5. Let \(f: \mathbf{X} \to \mathbf{Y}\) and \(f': \mathbf{X}' \to \mathbf{Y}'\) be morphisms of manifolds. Then \(f \times f': \mathbf{X} \times \mathbf{X}' \to \mathbf{Y} \times \mathbf{Y}'\) is a morphism. If \(a \in \mathbf{X}\) and \(a' \in \mathbf{X}'\) , we have:

\[
\begin{array}{l} \mathrm{T} _ {(a, a ^ {\prime})} (f \times f ^ {\prime}) = \mathrm{T} _ {a} (f) \times \mathrm{T} _ {a ^ {\prime}} (f ^ {\prime}) \\ \mathrm{rg} _ {(a, a ^ {\prime})} (f \times f ^ {\prime}) = \mathrm{rg} _ {a} f + \mathrm{rg} _ {a ^ {\prime}} f ^ {\prime}. \end{array}
\]

5.6.6. Let \(X_{1}\) , \(X_{2}\) and Z be three manifolds and f a morphism from \(X_{1} \times X_{2}\) to Z. Let \(a \in X_{1}\) and \(b \in X_{2}\) . The tangent linear map to the partial map \(x \mapsto f(x, b)\) (resp. \(y \mapsto f(a, y)\) ) of \(X_{1}\) (resp. \(X_{2}\) ) into Z is denoted \(T_{(a,b)}^{1}(f)\) (resp. \(T_{(a,b)}^{2}(f)\) ). If one identifies \(T_{(a,b)}(X_{1} \times X_{2})\) with \(T_{a}(X_{1}) \times T_{b}(X_{2})\) , we have:

\[
\mathrm{T} _ {(a, b)} (f). (u, v) = \mathrm{T} _ {(a, b)} ^ {1} (f). u + \mathrm{T} _ {(a, b)} ^ {2} (f). v
\]

for every \(u \in \mathrm{T}_a(\mathrm{X}_1)\) and every \(v \in \mathrm{T}_b(\mathrm{X}_2)\) .

5.6.7. ("Implicit function theorem"). With the hypotheses and notations of the preceding number, suppose moreover that \(\mathrm{T}_{(a,b)}^{2}(f)\) is bijective. Then there exists an open neighborhood U of a in \(X_{1}\) and an open neighborhood V of b in \(X_{2}\) with the following property: for every \(x \in U\) , there exists one and only one point \(g(x)\) of V such that \(f(x, g(x)) = f(a, b)\) , and the map g is a morphism from U to V. For every \(x \in U\) , one has

\[
\mathrm{T} _ {x} (g) = - \left(\mathrm{T} _ {(x, g (x))} ^ {2} (f)\right) ^ {- 1} \circ \mathrm{T} _ {(x, g (x))} ^ {1} (f).
\]

